\documentclass[11pt]{article}

\usepackage[letterpaper,margin=1in]{geometry}
\usepackage[T1]{fontenc}
\usepackage[utf8]{inputenc}
\usepackage{amsmath,amssymb}
\usepackage{amsthm}
\usepackage{booktabs}
\usepackage{enumitem}
\usepackage{float}
\usepackage{graphicx}
\usepackage{microtype}
\usepackage{xcolor}
\usepackage{xurl}
\usepackage[colorlinks=true,linkcolor=blue!55!black,urlcolor=blue!55!black,
            citecolor=blue!55!black]{hyperref}

\setlength{\parindent}{1.25em}
\setlength{\parskip}{0pt}
\setlist{itemsep=1pt,topsep=3pt,leftmargin=1.55em}
\newcommand{\code}[1]{\texttt{#1}}
\newcommand{\Bench}{\mathsf{B}}
\newtheorem{proposition}{Proposition}
\newtheorem{remark}{Remark}
\renewcommand{\qedsymbol}{$\blacksquare$}

\title{\vspace{-1.25cm}\textbf{Lexicographic Constraint Programming for\\
Recreational Softball Defensive Lineups: A Team Red Case Study}}
\author{David Russo\\\small Technical report}
\date{September 2026, revised September 11, 2026}

\begin{document}
\maketitle
\vspace{-0.75cm}

\begin{abstract}
We study the construction of seven-inning defensive lineups for recreational
slowpitch softball when the available roster exceeds the number of fielding
positions. The problem requires a legal assignment in every inning while
balancing playing time, respecting player-reported positional preferences, and
limiting disruptive field-position and bench transitions. We formulate the
problem as a constraint-programming model with a strict lexicographic
objective: playing-time inequity is minimized first, followed by non-explicit
assignments and several measures of lineup continuity. Eligibility distinguishes
explicitly selected positions from directionally inferred alternatives, while
pitcher and first base remain opt-in assignments. The model is implemented in
OR-Tools CP-SAT through bounded, status-aware phases suitable for interactive
use. For the fourteen-player Team Red snapshot studied here, the minimum
playing-time spread is three innings and the minimum scaled absolute deviation
is 56. A recorded legal incumbent fills all 70 defensive slots without a
non-explicit assignment, uses no more than two field positions per player, and
has continuity vector $(3,0,15,20,20)$. Its application status is
\emph{FEASIBLE}: the fairness and fallback values are certified, but complete
lexicographic optimality of the continuity vector is not claimed.
\end{abstract}

\textbf{Keywords:} sports scheduling; constraint programming; CP-SAT;
lexicographic optimization; playing-time fairness; human-centered operations

\section{Introduction}

Defensive scheduling in recreational slowpitch softball is a compact but
nontrivial assignment problem. A conventional defense contains ten positions:
pitcher (P), catcher (C), first base (1B), second base (2B), third base (3B),
shortstop (SS), left field (LF), left center field (LC), right center field
(RC), and right field (RF). A roster commonly contains more available players
than defensive positions. The resulting schedule must repeatedly decide who
fields, who occupies each position, and who is on the Bench.

The Open league considered here does not divide players by gender and does
not impose gender minima. That simplification removes one class of hard
constraints, but it does not make the practical objective obvious. Equal
playing time is important; explicit positional preferences should be honored;
players should not be moved unnecessarily; and a player who sits should, when
possible, sit in a coherent block rather than alternate between the field and
the Bench. These goals conflict. A smoother schedule is not necessarily a
more preference-respecting schedule, and a perfectly continuous schedule is
not necessarily fair.

The implementation resolves those conflicts by priority rather than by an
opaque weighted compromise. The order is substantive: legality precedes
fairness, fairness precedes aggregate non-explicit assignment count, and both
precede continuity. Lower-priority improvements are prohibited from purchasing
a worse higher-priority result. This formulation is particularly appropriate
for an interactive tool, where a user should not have to interpret a hidden
exchange rate between one inning of inequity and one position change.

This work makes three contributions:
\begin{enumerate}
    \item it formulates recreational defensive scheduling as a finite-horizon
    state-assignment problem with an explicit bench state;
    \item it separates preference declarations from directionally inferred
    eligibility and orders fairness, non-explicit assignment, and continuity
    objectives lexicographically; and
    \item it develops a bounded CP-SAT procedure that preserves certified
    higher-priority values while returning a legal incumbent under an
    interactive time budget.
\end{enumerate}

The presentation follows standard assignment-modeling practice
\cite{bertsimas,williams}, but the repeated within-game assignment, directed
eligibility expansion, and explicit bench state produce an application-specific
multiobjective structure. Mathematical propositions establish structural and
instance-specific bounds; solver statuses certify only those objective tiers
proved optimal; regression tests check implementation conformance; and timed
experiments characterize, rather than guarantee, latency and incumbent quality.

\section{Operational Assumptions}

The model optimizes one game at a time under the following assumptions.
\begin{itemize}
    \item A regulation planning horizon contains seven innings.
    \item Availability is constant throughout the game. Late arrival and
    early departure are outside the current input model.
    \item Between eight and fifteen players may be selected as available.
    A stored roster may be larger, but a single solve is capped at fifteen.
    \item Every active position is filled in every modeled inning, and one
    player cannot occupy two simultaneous positions.
    \item Gender has no legal or objective effect in the Open profile.
    Team Red therefore stores an internal \code{Unspecified} sentinel and
    exposes no gender control.
    \item A checked position is an unranked explicit preference. Derived
    positions are legal fallbacks, not additional preferences.
    \item Every player may catch. Pitcher and first base require explicit
    opt-in because neither is inferred from another position.
    \item The Bench state means only that a player is not fielding in that
    inning. It does not remove the player from the batting order or model an
    offensive substitution.
    \item For a given confirmed input snapshot, the returned schedule is
    authoritative. The designated lineup caller may alter the inputs and solve
    again, but no approval gate or hand edit modifies the output.
\end{itemize}

Let $N$ denote the number of available players. The active position set
$\mathcal S(N)$ and its cardinality $K(N)=|\mathcal S(N)|$ are fixed before
optimization according to Table~\ref{tab:active-positions}.

\begin{table}[htbp]
\centering
\caption{Active defensive positions under the Open rules profile.}
\label{tab:active-positions}
\begin{tabular}{ccl}
\toprule
Available players $N$ & $K(N)$ & $\mathcal S(N)$\\
\midrule
$10\leq N\leq15$ & 10 & P, C, 1B, 2B, 3B, SS, LF, LC, RC, RF\\
$N=9$ & 9 & all standard positions except RF\\
$N=8$ & 8 & all standard positions except C and RF\\
\bottomrule
\end{tabular}
\end{table}

An input with fewer than eight or more than fifteen available players is
rejected before the constraint model is constructed. If ten or more players
are available, the optimizer does not silently reduce the defense to nine to
work around insufficient eligibility; it reports the eligibility failure.

\section{Preference and Eligibility Semantics}

Let the complete position universe be
\[
\mathcal{S}=\{\mathrm{P,C,1B,2B,3B,SS,LF,LC,RC,RF}\},
\]
and let $\varnothing\neq Q_p\subseteq\mathcal S$ be player $p$'s explicitly
selected positions. Define the direct eligibility-expansion map $H_f$ by
\[
\begin{aligned}
H_f(\mathrm{SS})&=\{\mathrm{3B,2B}\}, &
H_f(\mathrm{3B})&=\{\mathrm{2B}\},\\
H_f(\mathrm{LC})&=\{\mathrm{LF,RC,RF}\}, &
H_f(\mathrm{LF})&=\{\mathrm{RC,RF}\},\\
H_f(\mathrm{RC})&=\{\mathrm{RF}\}, &
H_f(s)&=\varnothing\quad\text{otherwise}.
\end{aligned}
\]
Thus selecting SS, for example, implies eligibility at 3B and 2B; the notation
denotes eligibility expansion, not preference ranking or dominance. The map is
applied directly rather than transitively:
$h(Q_p)=\bigcup_{s\in Q_p}H_f(s)$.
Every eligible non-explicit assignment---whether produced by $H_f$, universal
catcher eligibility, or the eight-player rule below---has unit fallback cost.
The lexicographic objective fixes the optimal or best-found aggregate fallback
value before continuity is optimized. Among schedules having that same value,
continuity may choose RF for an LC-preferring player or redistribute which
player receives a fallback. The model does not separately optimize the maximum
or distribution of individual fallback burden.

The eligible active set is
\begin{equation}
E_p(N) = \bigl(Q_p\cup h(Q_p)\cup\{\mathrm{C}\}\cup R_p(N)\bigr)
      \cap\mathcal{S}(N),
\label{eq:eligibility}
\end{equation}
where
\[
R_p(N) =
\begin{cases}
\{\mathrm{RC}\}, & N=8\text{ and }\mathrm{RF}\in Q_p,\\
\varnothing, & \text{otherwise.}
\end{cases}
\]
The special RF-to-RC rule applies only to the exactly eight-player defense,
when both C and RF are inactive. It does not apply at nine players. Universal
catcher eligibility is operational policy, while P and 1B are absent from
$h(\cdot)$ and must be explicitly selected.

The fallback set is
\[
F_p(N)=E_p(N)\setminus Q_p.
\]
The model may assign $p$ to $F_p$ to preserve feasibility or fairness, but it
minimizes such innings before evaluating any continuity term. Thus eligibility
is broader than preference without erasing the semantic distinction between
the two.

\section{Constraint-Programming Formulation}

For a fixed accepted roster size $N$, write $\mathcal S_N=\mathcal S(N)$,
$K=K(N)$, $E_p=E_p(N)$, and $F_p=F_p(N)$. Table~\ref{tab:notation}
collects the principal notation.

\begin{table}[htbp]
\centering
\caption{Principal sets, variables, and aggregate quantities.}
\label{tab:notation}
\small
\begin{tabular}{ll}
\toprule
Symbol & Definition\\
\midrule
$\mathcal P,\mathcal I,\mathcal S_N$ & players, innings, and active positions\\
$Q_p,E_p,F_p$ & explicit, eligible, and eligible non-explicit positions for $p$\\
$x_{ips},b_{ip}$ & field-assignment and bench-state indicators\\
$t_p,d_p,\Delta$ & innings played, scaled deviation, and playing-time spread\\
$u_{ps},v_p,e_p$ & position-use indicator, distinct positions, and excess above two\\
$f$ & total eligible non-explicit assignment innings\\
$o_1,o_2$ & numbers of exact one- and two-inning state runs\\
$v,e,a$ & aggregate distinct positions, excess positions, and transitions\\
$H=7K$ & total fielding slots\\
\bottomrule
\end{tabular}
\end{table}

\subsection{Sets and decision variables}

Let $\mathcal{P}=\{1,\ldots,N\}$ be the available players and
$\mathcal{I}=\{1,\ldots,7\}$ the innings. The primary binary variable is
\[
x_{ips}=
\begin{cases}
1,&\text{if player $p$ fields at position $s$ in inning $i$},\\
0,&\text{otherwise,}
\end{cases}
\]
for $i\in\mathcal{I}$, $p\in\mathcal{P}$, and $s\in E_p$. A second binary
variable $b_{ip}$ equals one exactly when player $p$ is on the Bench in inning
$i$. Omitting $x$ variables for positions outside $E_p$ makes an ineligible
fielding assignment structurally impossible.

All variables denoted by $x$, $b$, $w$, $a$, $q^{(1)}$, $q^{(2)}$, and $u$
are binary. The count variables satisfy $t_p\in\{0,\ldots,7\}$,
$v_p\in\{0,\ldots,K\}$, and
$e_p\in\{0,\ldots,\max(0,K-2)\}$.

Each active position is filled exactly once:
\begin{equation}
\sum_{p:s\in E_p}x_{ips}=1
\qquad\forall i\in\mathcal{I},\ s\in\mathcal{S}_N.
\label{eq:position-coverage}
\end{equation}
Each player occupies exactly one field-or-Bench state:
\begin{equation}
b_{ip}+\sum_{s\in E_p}x_{ips}=1
\qquad\forall i\in\mathcal{I},\ p\in\mathcal{P}.
\label{eq:player-state}
\end{equation}
Constraint~\eqref{eq:player-state} simultaneously prevents double assignment and explicitly
represents the Bench. Because the Open profile has no gender rule, the model
contains no gender-count or gender-placement constraint.

Before these variables are constructed, input validation rejects blank or
duplicate names, empty or invalid preference sets, and out-of-range player
counts. A distinct lineup preflight then considers the bipartite graph
$G=(\mathcal S_N\cup\mathcal P,A)$, where $(s,p)\in A$ if and only if
$s\in E_p$.

\begin{proposition}[One-inning matching criterion]
\label{prop:matching}
Under the Open profile, a legal seven-inning schedule satisfying only the hard
assignment constraints exists if and only if $G$ has a matching that covers
$\mathcal S_N$. Equivalently, for every $T\subseteq\mathcal S_N$,
\[
\left|\bigcup_{s\in T}\{p\in\mathcal P:s\in E_p\}\right|\ge |T|.
\]
\end{proposition}
\begin{proof}
Constraints~\eqref{eq:position-coverage}--\eqref{eq:player-state} restricted to
one inning are precisely a matching covering $\mathcal S_N$; Hall's theorem
gives the stated equivalence. Availability and eligibility are invariant across
innings, and the model has no cross-inning hard constraint. Repeating a covering
matching therefore constructs a legal seven-inning schedule. Conversely, the
assignment in any one inning of a legal schedule supplies such a matching.
\end{proof}

The implementation first reports any uncovered singleton position. Otherwise,
it reports a deterministic minimum-cardinality Hall-deficient subset when no
covering matching exists.

\subsection{Playing-time quantities}

The number of fielding innings received by player $p$ is
\[
t_p=\sum_{i\in\mathcal{I}}\sum_{s\in E_p}x_{ips}.
\]
Let
\[
t^{\max}=\max_{p\in\mathcal{P}}t_p,
\qquad
t^{\min}=\min_{p\in\mathcal{P}}t_p,
\qquad
\Delta=t^{\max}-t^{\min}.
\]
There are $H=7K$ fielding slots. To avoid fractional expressions, define the
scaled deviation from an equal share as
\[
d_p=\left|Nt_p-H\right|.
\]
Minimizing $\sum_p d_p$ is equivalent to minimizing total absolute deviation
from $H/N$, up to the positive factor $N$.

\subsection{State runs and transitions}

For continuity purposes, each player has the state set
\[
\mathcal{R}_p=E_p\cup\{\Bench\}.
\]
Let $w_{ipr}$ equal $x_{ipr}$ for a field position $r\in E_p$ and equal
$b_{ip}$ for $r=\Bench$. A run is a maximal consecutive subsequence of
identical states. Consequently,
\[
\mathrm{SS}\rightarrow\Bench\rightarrow\mathrm{SS}
\]
contains three runs and two transitions; it is not interpreted as one
interrupted shortstop assignment.

For $i>1$, a transition-start variable $a_{ipr}$ represents the conjunction
$w_{ipr}\land\neg w_{i-1,p,r}$ through
\[
a_{ipr}\leq w_{ipr},\qquad
a_{ipr}+w_{i-1,p,r}\leq1,\qquad
a_{ipr}\geq w_{ipr}-w_{i-1,p,r}.
\]
Exactly one new state starts whenever adjacent innings have different states.
The first inning is not charged as a transition.

With boundary values $w_{0pr}=w_{8pr}=0$, the binary one-inning-run indicator
$q^{(1)}_{ipr}$, for $i\in\{1,\ldots,7\}$, recognizes the pattern $0,1,0$:
\[
\begin{aligned}
q^{(1)}_{ipr}&\leq w_{ipr},\\
q^{(1)}_{ipr}&\leq 1-w_{i-1,p,r},\\
q^{(1)}_{ipr}&\leq 1-w_{i+1,p,r},\\
q^{(1)}_{ipr}&\geq
w_{ipr}-w_{i-1,p,r}-w_{i+1,p,r}.
\end{aligned}
\]
For $i\in\{1,\ldots,6\}$, the binary two-inning-run indicator
$q^{(2)}_{ipr}$ recognizes the exact pattern $0,1,1,0$ through
\[
\begin{aligned}
q^{(2)}_{ipr}&\leq w_{ipr}, &
q^{(2)}_{ipr}&\leq w_{i+1,p,r},\\
q^{(2)}_{ipr}&\leq 1-w_{i-1,p,r}, &
q^{(2)}_{ipr}&\leq 1-w_{i+2,p,r},\\
q^{(2)}_{ipr}&\geq w_{ipr}+w_{i+1,p,r}
                -w_{i-1,p,r}-w_{i+2,p,r}-1.
\end{aligned}
\]
Binary domains make each linearization an if-and-only-if characterization.
Both definitions apply to field positions and $\Bench$. They measure rather
than prohibit short runs; feasibility and higher-priority fairness remain
paramount.

\subsection{Positional variety}

Let $u_{ps}$ equal one if player $p$ uses field position $s$ at least once.
The standard linkage is
\[
x_{ips}\leq u_{ps}\quad
\forall p\in\mathcal P,\ s\in E_p,\ i\in\mathcal I,
\qquad
u_{ps}\leq\sum_{i\in\mathcal I}x_{ips}\quad
\forall p\in\mathcal P,\ s\in E_p.
\]
Define the integer number of distinct field positions as
\[
v_p=\sum_{s\in E_p}u_{ps}.
\]
Mathematically, $e_p=\max\{v_p-2,0\}$. The implementation gives the
nonnegative bounded integer $e_p$ the constraints
\[
e_p\geq v_p-2,\qquad e_p\geq0.
\]
Because its objective coefficient is positive, the lower bound is tight in any
minimizer of the continuity objective; equality is not imposed as a primitive
constraint. An earlier incumbent retained after a time-limited phase need not
carry a certificate of this tightness.
Bench does not contribute to $v_p$ or $e_p$. The two-position target is soft:
the model may exceed it to protect legality, fairness, or the fallback count.

\section{Lexicographic Objective}

Define
\[
\begin{aligned}
f&=\sum_{i\in\mathcal I}\sum_{p\in\mathcal P}
     \sum_{s\in F_p}x_{ips},\\
o_1&=\sum_{p\in\mathcal P}\sum_{r\in\mathcal R_p}
     \sum_{i=1}^{7}q^{(1)}_{ipr}, &
o_2&=\sum_{p\in\mathcal P}\sum_{r\in\mathcal R_p}
     \sum_{i=1}^{6}q^{(2)}_{ipr},\\
e&=\sum_{p\in\mathcal P}e_p, &
v&=\sum_{p\in\mathcal P}v_p,\\
a&=\sum_{p\in\mathcal P}\sum_{r\in\mathcal R_p}
     \sum_{i=2}^{7}a_{ipr}.
\end{aligned}
\]
Let $\boldsymbol y$ collect all decision variables, and let $\mathcal X$ be the
feasible set defined by the preceding constraints. The ideal mathematical
problem is
\begin{equation}
\underset{\boldsymbol y\in\mathcal X}{\operatorname{lexmin}}\;
z(\boldsymbol y)=\left(\Delta,\sum_{p\in\mathcal P}d_p,
f,o_1,e,o_2,v,a\right).
\label{eq:lex-objective}
\end{equation}
Thus fairness has two ordered components, followed by the aggregate fallback
count and five ordered continuity components.

The implementation scalarizes only within the fairness and continuity blocks.
Since $0\leq t_p\leq7$, a valid bound is
$D^{UB}=N\max\{H,|7N-H|\}$. The fairness objective is
\begin{equation}
J_F=(D^{UB}+1)\Delta+\sum_p d_p.
\label{eq:fairness-objective}
\end{equation}
The multiplier ensures that any one-unit reduction in spread dominates every
possible change in total deviation.

After the fairness pair and fallback value have been fixed, the continuity
objective is
\begin{equation}
J_C=c_1o_1+c_e e+c_2o_2+c_vv+a,
\label{eq:continuity-objective}
\end{equation}
using the valid bounds
\[
a^{UB}=6N,\quad v^{UB}=NK,\quad o_2^{UB}=6N,
\quad e^{UB}=N\max\{0,K-2\}
\]
and the recurrence
\[
\begin{aligned}
c_v&=a^{UB}+1,\\
c_2&=c_vv^{UB}+a^{UB}+1,\\
c_e&=c_2o_2^{UB}+c_vv^{UB}+a^{UB}+1,\\
c_1&=c_ee^{UB}+c_2o_2^{UB}+c_vv^{UB}+a^{UB}+1.
\end{aligned}
\]

\begin{proposition}[Lexicographic equivalence]
\label{prop:lex-equivalence}
Minimizing $J_C$ is equivalent to lexicographically minimizing
$(o_1,e,o_2,v,a)$ over any feasible slice on which the fairness pair and
fallback value are fixed.
\end{proposition}
\begin{proof}
Each coefficient exceeds the largest possible contribution of all terms to its
right. Consequently, a unit improvement in the first differing component
strictly decreases $J_C$, irrespective of every lower-priority component; the
converse follows by applying the same domination argument to the first differing
component of two feasible vectors.
\end{proof}

It follows that continuity cannot increase the fixed fallback value, although
it may redistribute the same number of fallback assignments among players. A
$\Bench$ run is evaluated by the same logic as a field-position run. Exact
two-inning runs are minimized only after one-inning runs and excess positions;
the model therefore never purchases a reduction in $o_2$ by worsening $o_1$ or
$e$.

\section{Lexicographic Solution Method}

The model is implemented with the OR-Tools CP-SAT solver, which supports
Boolean and integer constraint models and distinguishes statuses including
\emph{OPTIMAL}, \emph{FEASIBLE}, and \emph{INFEASIBLE} \cite{ortoolscpsat}.
The interactive default is a nominal five-second multi-phase solve target measured
only after model construction. It is not a hard end-to-end wall-clock bound:
model-building time, Python and Streamlit overhead, CP-SAT time-limit behavior,
and the implementation's minimum per-phase and subphase allocations of
0.05--0.1 seconds can make observed elapsed time exceed five seconds. The implementation uses the
following nominally bounded phases:
\begin{enumerate}
    \item \textbf{Balanced-count feasibility test.} Write $H=qN+r$, where
    $0\le r<N$. Without eligibility restrictions, the balanced count multiset
    has $N-r$ values of $q$ and $r$ values of $q+1$, with
    $\Delta=\mathbf 1_{\{r>0\}}$ and $\sum_p d_p=2r(N-r)$. Add
    $\lfloor H/N\rfloor\leq t_p\leq\lceil H/N\rceil$ for every player and
    test feasibility. A feasible assignment certifies that both arithmetic
    lower bounds remain attainable under all hard constraints.
    \item \textbf{General fairness search.} If that slice cannot be established,
    minimize $J_F$ within the designated fairness budget.
    \item \textbf{Preference phase.} Fix the achieved spread and deviation,
    minimize $f$, and fix the exact fallback incumbent before any continuity
    phase. If optimality is not proved, the incumbent is still frozen rather
    than silently relaxed.
    \item \textbf{One-inning phase.} For at most 0.25 seconds, test the
    lower-bound slice in which the only single-inning field run is forced by
    $t_p=1$ and the only single-inning $\Bench$ run is forced by $t_p=6$.
    Within this hard slice, minimize $e$ as look-ahead to the next tier. If
    feasibility is not established within that time slice,
    minimize $o_1(E^{UB}+1)+e$ for up to 2.5 seconds, where $E^{UB}$ bounds total excess positions,
    and preserve the one-inning incumbent. If that search produces no
    candidate, retain the already legal fairness/fallback incumbent and
    continue with the complete continuity objective. The coefficient keeps
    every unit of $o_1$ strictly more important than all possible changes in $e$.
    \item \textbf{Two-position phase.} If the preceding $o_1$ tier is proved,
    test the $e=0$ slice inside the existing phase budget and otherwise
    minimize $e$ with the unspent remainder. A proved value is fixed. If
    $o_1$ remains unproved, do not independently cap $e$: retain the complete
    ceiling $J_C\leq J_C^{inc}$, where $J_C^{inc}$ is the accepted incumbent's
    integer objective value, so a later improvement in $o_1$ may
    legitimately spend lower-priority quality.
    \item \textbf{Two-inning phase.} When $o_1$ and $e$ are both proved,
    minimize $o_2c_2+vc_v+a$; otherwise keep the complete $J_C$ objective
    active. Run two concurrent four-worker searches sharing the existing
    eight-worker computational envelope. One is an ordinary proof-capable
    CP-SAT run; the second is a hint-seeded large-neighborhood search with a
    different seed. Retain the lexicographically better legal candidate. Fix
    $o_2$ only when its prefix is proved; otherwise preserve the full weighted
    incumbent ceiling.
    \item \textbf{Tail-continuity phase.} Use the remaining nominal budget to
    minimize $vc_v+a$ only when every earlier continuity tier is proved; if
    not, continue minimizing $J_C$. Because a hint is not a quality guarantee,
    constrain the complete weighted objective to be no larger than the
    incumbent and accept a timed candidate only when its full
    $(o_1,e,o_2,v,a)$ vector is lexicographically no worse than that incumbent.
\end{enumerate}

An \emph{UNKNOWN} status from a time-limited slice is not interpreted as
infeasibility. Every candidate is required to have a solved status and is then
compared with the accepted incumbent by tuple lexicographic order, not
componentwise order. A proved higher-priority value is fixed before the next
phase begins; without such a proof, the complete incumbent ceiling remains
active.

The application synthesizes its public status from all phase results. It reports
\emph{OPTIMAL} if and only if every prefix proof flag is true and the final tail
solve returns CP-SAT \emph{OPTIMAL}. Otherwise any returned legal incumbent is
labeled \emph{FEASIBLE}, even when one or more prefix tiers are proved. The
latter label denotes absence of a complete lexicographic optimality certificate,
not uncertainty about legality.

\section{Computational Evidence}

Here For The Beer (HFTB) is a separate fifteen-player co-ed benchmark instance;
it is not Team Red evidence. On the recorded macOS reference environment at
commit \code{fa92000}, twenty measured HFTB calls to \code{optimize\_game}
were distributed across four fresh Python processes after one excluded warm-up
per process. The reported duration includes model construction and all solver
phases, but excludes process startup, roster loading, and the excluded warm-up.
All twenty preserved the proved fairness/fallback prefix
$(4,210,0)$. Nineteen met the timed continuity target
$(5,0,9,20,21)$ lexicographically: twelve returned $(5,0,9,19,20)$ and seven
returned $(5,0,9,20,20)$. One legal incumbent returned
$(5,0,10,20,21)$ and is retained explicitly in the evidence rather than hidden
by an aggregate. Call duration had p95 5.062 seconds and maximum 5.063 seconds;
the fairness/fallback prefix had p95 0.138 seconds and maximum 0.151 seconds.
Every result remained \emph{FEASIBLE}. The configured experimental gate was an
acceptance proportion of at least 0.90; the observed empirical proportion was
$19/20=0.95$. No population guarantee or confidence bound is claimed. The p95
statistic is the nearest-rank order statistic for $n=20$, not an estimate of a
service-level guarantee.

\begin{table}[htbp]
\centering
\caption{Recorded HFTB benchmark results at revision \code{fa92000}.}
\label{tab:hftb-results}
\small
\begin{tabular}{lr}
\toprule
Quantity & Recorded value\\
\midrule
Measured repetitions & 20\\
Certified fairness/fallback prefix & $(4,210,0)$ in 20/20 runs\\
Continuity vector $(5,0,9,19,20)$ & 12 runs\\
Continuity vector $(5,0,9,20,20)$ & 7 runs\\
Continuity vector $(5,0,10,20,21)$ & 1 run\\
Target-satisfying proportion & $19/20=0.95$\\
Call duration, p95 / maximum & 5.062 / 5.063 seconds\\
Prefix duration, p95 / maximum & 0.138 / 0.151 seconds\\
\bottomrule
\end{tabular}
\end{table}

\section{Team Red Case Study}

This case study uses the fourteen-player Team Red snapshot reproduced in
Table~\ref{tab:team-red-roster}, with every player initially available under
locked Open rules.

\begin{table}[htbp]
\centering
\caption{Team Red roster snapshot and explicit positional preferences.}
\label{tab:team-red-roster}
\begin{tabular}{ll@{\qquad}ll}
\toprule
Player & Explicit preferences & Player & Explicit preferences\\
\midrule
David R & 2B, 3B, RC & Justin & 1B, RC\\
Andrew & 1B, 2B & Kevin & 2B, LC\\
David & RF & Dung & P\\
Heather & LF, LC & Ryan & 1B, 3B\\
Tristan & C, 3B & Brad & C, LF\\
AP & SS, LF & Chris & C, 2B, RF\\
Marty & SS, LC & Shaz & C, RF\\
\bottomrule
\end{tabular}
\end{table}

\begin{table}[H]
\centering
\caption{Representative Team Red defensive schedule.}
\label{tab:team-red-lineup}
\scriptsize
\resizebox{\textwidth}{!}{%
\begin{tabular}{c@{\quad}llllllllll}
\toprule
Inning & P & C & 1B & 2B & 3B & SS & LF & LC & RC & RF\\
\midrule
1 & Dung & Brad & Andrew & Chris & Ryan & Marty & Heather & Kevin & David R & David\\
2 & Dung & Brad & Justin & Chris & Ryan & Marty & Heather & Kevin & David R & David\\
3 & Dung & Shaz & Justin & Chris & Tristan & AP & Heather & Kevin & David R & David\\
4 & Dung & Shaz & Justin & Andrew & Tristan & AP & Heather & Kevin & David R & David\\
5 & Dung & Shaz & Ryan & Andrew & Tristan & AP & Brad & Marty & David R & David\\
6 & Dung & Shaz & Ryan & Andrew & Tristan & AP & Brad & Marty & Justin & Chris\\
7 & Dung & Shaz & Ryan & Andrew & Tristan & AP & Brad & Marty & Justin & Chris\\
\bottomrule
\end{tabular}}
\end{table}

Table~\ref{tab:team-red-lineup} records one machine-checked incumbent produced
by the nominal five-second procedure. It is included as a primal feasibility
witness, not as a claim of uniqueness.

\begin{proposition}[Optimal Team Red fairness profile]
\label{prop:team-red-fairness}
For the snapshot in Table~\ref{tab:team-red-roster}, every feasible schedule has
$\Delta\ge3$. Conditional on $\Delta=3$, the minimum scaled absolute deviation
is 56, attained by the inning-count multiset
$\{7^1,5^{11},4^2\}$.
\end{proposition}
\begin{proof}
Here $N=14$, $K=10$, and $H=70$. Pitcher is active in every inning, and Dung is
the unique player whose eligible set contains P. Let $p_D$ denote Dung. Hence
$x_{ip_D\mathrm P}=1$ for every $i\in\mathcal I$ and $t_{p_D}=7$.
If $\Delta\le2$, every other player must receive at least five innings, implying
$\sum_pt_p\ge7+13(5)=72>H$, a contradiction. Thus $\Delta\ge3$.

The schedule in Table~\ref{tab:team-red-lineup} has count multiset
$\{7^1,5^{11},4^2\}$ and therefore establishes $\Delta\le3$. For the second
tier, let $a_j$ count players with $j$ innings in a spread-three schedule.
Because $4\le t_p\le7$,
$a_4+a_5+a_6+a_7=14$ and
$4a_4+5a_5+6a_6+7a_7=70$, whence $a_4=a_6+2a_7$.
Moreover $a_7\ge1$, and
\[
\sum_pd_p=14a_4+14a_6+28a_7=28a_6+56a_7\ge56.
\]
The displayed schedule attains 56, completing both bounds.
\end{proof}

For the complete vector in~\eqref{eq:lex-objective}, the displayed incumbent has
\[
z=(3,56,0,3,0,15,20,20).
\]
The zero fallback count is optimal conditional on the certified fairness pair
because $f\ge0$. The application status was \emph{FEASIBLE}, so no optimality
claim is made for $(o_1,e,o_2,v,a)=(3,0,15,20,20)$. Multiple name-by-name
schedules may have the same vector.

The recorded run used code revision \code{ff382de}, roster SHA-256
\nolinkurl{4b3f2715e9dbb40dc8dfac945ea060953f6c31b17fd6d8be2d94af0ff33b9ad3},
Python 3.9.6, OR-Tools 9.15.6755, and macOS 26.6.2 on arm64. The measured
\code{optimize\_game} call took 5.073 seconds under the nominal five-second
budget. This single recorded execution is a reproducible artifact datum, not a
sampling statement about latency or incumbent quality.

The instance also exposes availability risk clearly. If Dung is marked
unavailable, no remaining player has opted into P. Simultaneous absences by
Justin, Andrew, and Ryan leave eleven players but remove all 1B coverage, while
simultaneous absences by AP and Marty leave twelve players but remove all SS
coverage. Adequate headcount therefore does not guarantee positional
feasibility. The correct response is an infeasibility message identifying the
uncovered position, not an invented assignment and not a reduced nine-player
lineup. The absence of any catcher opt-in, by itself, cannot leave C uncovered
because $\mathrm C\in E_p$ for every player; another Hall conflict could still
make the lineup infeasible.

\section{Verification and Falsification Strategy}

Validation is organized around properties rather than visual plausibility.
At prepublication revision \code{ff382de}, a local run recorded 144 passing
pytest cases. The Playwright suite defines eight scenario functions, each
configured for a mobile-touch project and a desktop-mouse project, yielding 16
project-test executions. These are implementation-conformance checks; they are
not substitutes for the analytic bounds or solver certificates stated above.
The principal checks include:
\begin{enumerate}
    \item \textbf{Legality invariants:} seven innings, exact active-position
    coverage, no duplicated player within an inning, valid eligibility, and
    reported summary counts recomputed from returned assignments, with separate
    assertions for eligibility.
    \item \textbf{Open-profile boundaries:} legal eight-, nine-, and ten-
    fielder position sets; no gender error in Open; no nine-player RF-to-RC
    special promotion; universal C fallback; and explicit P/1B requirements.
    \item \textbf{Team Red contract:} exact names and preferences, internal
    \code{Unspecified} gender, locked Open rules, all players initially
    available, fair inning distribution, and no player assigned to more than
    two field positions.
    \item \textbf{Preference-priority witness:} the fixture implemented by
    \path{test_explicit_preferences_beat_a_strictly_smoother_fallback_schedule}
    in \code{tests/test\_optimizer.py} has a
    solver-proven seven-fallback optimum with continuity vector
    $(5,1,16,20,21)$. Promoting Third Specialist's already-eligible 2B from a
    fallback to an explicit preference leaves the eligibility sets unchanged
    and yields the smoother vector $(5,0,14,20,20)$. Rescored under the original
    preferences, that exhibited schedule uses nine fallbacks. The five vector
    components are, respectively, one-inning runs, excess positions,
    two-inning runs, total distinct positions, and adjacent transitions; the
    first improvement is excess positions, from one to zero. The original
    optimizer nevertheless selects the seven-fallback result. The test detects
    violations on this fixture; the general non-worsening property follows from
    imposing the exact constraint $f=f^{inc}$ before continuity search.
    \item \textbf{Mutation test:} weakening the exact fallback-incumbent freeze
    to a one-sided inequality is detected. The regression inspects the model
    domain and requires the exact singleton interval $[f,f]$.
    \item \textbf{Interaction tests:} real individual position taps, complete
    name entry, gender change where applicable, save/reopen persistence,
    optimization, the complete semantic seven-inning matrix, CSV download, invalid
    optimization recovery, removal, reset, and independent browser sessions.
\end{enumerate}

The strict preference witness is important because a test that merely observes
zero fallback use is vacuous: continuity cannot demonstrate an illegal trade
when no trade is available. The constructed counterexample supplies a smoother
but preference-worse alternative and verifies that the implemented priority
rejects it. This is closer to falsification than to screenshot confirmation.

\section{Human and Operational Interpretation}

The objective order encodes several practical slowpitch considerations.
Continuous field-position blocks reduce repeated instruction, warm-up changes,
and uncertainty between innings only as an intended operational proxy; no user
study estimates those causal effects. Modeling $\Bench$ as a state prevents the
objective from ignoring repeated movement between the field and bench.
Consecutive bench innings are likewise intended to simplify communication from
a phone in a dugout, but that usability claim has not been experimentally
evaluated.

The hierarchy supplies controlled resilience without pretending all positions
are interchangeable. Its SS-to-3B/2B, 3B-to-2B, and LC-to-LF/RC/RF pathways
are declared team-policy eligibility rules, not universal claims that one
position is easier than another. The separate fallback phase prevents
continuity from increasing the fixed aggregate fallback value, while still
permitting redistribution among schedules tied at that value. Pitching and
first base remain explicit-only choices.

Operational authority is scoped to an input snapshot, not to every open web
session. Team Red's defaults mark all fourteen players available, but that is
not attendance confirmation. Before solving, one designated lineup caller must
mark every scratch. Any later availability or preference change requires a new
solve, and the newly displayed or downloaded matrix supersedes the previous
one without hand editing. Because browser sessions are independent and not
synchronized, the team should distribute the result from that single source,
preferably by downloading the final matrix before the first pitch.

No captain-review stage follows optimization. Such a stage would blur the
contract by presenting a computed schedule and then asking a human to repair
it without rechecking legality, fairness, or continuity. In the implemented
workflow, changed judgments are expressed as changed availability or
preferences and submitted to a new solve. The optimizer, rather than an
unvalidated hand edit, remains the source of the displayed and downloaded
lineup.

\section{Limitations and Threats to Validity}

The model is intentionally narrow. Bench means not fielding; it does not alter
batting order. The model does not optimize offensive substitutions, opponent
tendencies, score-dependent tactics, player fatigue, or inning-specific
availability. It balances playing time within one game and does not carry
season-to-date equity between sessions. A shortened game simply leaves later
rows unused and may make realized playing time less equal than the planned
seven-inning distribution; extra innings are outside the fixed horizon. The
hierarchy is a declared operational policy, not an empirical estimate of every
individual's defensive ability.

Universal catcher eligibility is likewise a policy assumption. The software
correctly applies it to every roster and does not offer a catcher opt-out.
If a future league or team rejects that assumption, the input schema and hard
eligibility rules would need an explicit change; silently weakening the rule
inside the solver would be inappropriate.

The objective treats all explicit preferences as equal and all eligible
non-explicit assignments as equal. It does not know that one player strongly
favors LC over LF unless the roster is edited to reflect only the intended
choices. A future
ranked-preference model could add tiers, but it would require additional input
and a revised priority contract.

Finally, each CP-SAT phase is assigned a nominal time limit. A
\emph{FEASIBLE} result is guaranteed to satisfy the encoded hard constraints
but may lack a proof that every reported objective tier is globally best.
Conservative status reporting makes this uncertainty visible, and the web
interface reports observed elapsed time. This paper does not infer latency
percentiles or a hard wall-clock guarantee from the mere existence of a
benchmark utility. Any expansion beyond the eight-to-fifteen-player scope or
seven-inning horizon requires renewed, reproducible timing measurement.

\section{Conclusion}

Open-division defensive scheduling is best viewed as a lexicographic state-
assignment problem. The model fills every required position, equalizes playing
time as far as eligibility permits, preserves explicit preferences before
using controlled eligibility expansion, and then optimizes ordered continuity
measures over field and $\Bench$ states. Fixing the aggregate fallback value
before continuity search preserves the distinction between explicitly declared
positions and merely eligible alternatives, although individual fallback
assignments may still differ among aggregate ties.

For Team Red, Proposition~\ref{prop:team-red-fairness} proves a minimum playing-time spread of three and a
minimum scaled deviation of 56. The recorded incumbent uses no fallback
assignments, assigns no player to more than two field positions, and has
continuity vector $(3,0,15,20,20)$; the public \emph{FEASIBLE} status expressly
withholds a claim of complete continuity optimality. The staged procedure is
therefore suitable as a legal-incumbent generator on an interactive time scale,
subject to the empirical and modeling limitations stated above. Natural
extensions include inning-specific availability, season-level equity, ranked
preferences, and empirical calibration of player--position suitability.

\appendix
\section{Algorithmic Summary}

The implemented procedure may be summarized as follows:
\begin{enumerate}
    \item normalize and validate the selected available players;
    \item select the Open active-position set from $N$;
    \item derive $E_p$ and $F_p$ while preserving $Q_p$;
    \item preflight a distinct-player matching that covers all active positions;
    \item construct field, $\Bench$, fairness, run, position-use, and transition
    variables;
    \item attempt the balanced-count feasibility certificate, otherwise optimize
    fairness;
    \item constrain the achieved fairness pair exactly;
    \item minimize and freeze fallback innings;
    \item minimize and preserve one-inning state runs;
    \item minimize excess positions and prevent subsequent worsening;
    \item minimize the lexicographically weighted continuity expression;
    \item translate the legal incumbent into an inning matrix and player
    summary; and
    \item label the result \emph{OPTIMAL} only when all required proofs exist,
    otherwise label the legal result \emph{FEASIBLE}.
\end{enumerate}

\section{Application and Reproducibility Artifacts}

The public neutral application is available at
\url{https://softball-fielding-optimizer.streamlit.app/}. It presents an
explicit first-load choice among Here For The Beer, Team Red, and a blank
roster. Team Red loads the roster in \code{team\_red\_roster.csv} under locked
Open rules. The existing Here For The Beer deployment remains available at
\url{https://here-for-the-beer.streamlit.app/} and retains its Co-ed defaults.

The implementation, normative specification, test suite, and benchmark
utilities are maintained in the public repository at
\url{https://github.com/davidluizrusso/softball-fielding}. Streamlit Community
Cloud deploys an application from a GitHub repository and entrypoint and
provides a shareable \code{streamlit.app} URL \cite{streamlitdeploy}. Python
dependencies are declared in \code{requirements.txt}, consistent with the
platform's dependency mechanism \cite{streamlitdeps}.

During a solve, the web interface reports the active phase and displays a small
animated stick figure. This cue is intended to distinguish active computation
from an unregistered input event; no usability experiment is claimed.

The command-line verification entry points are:
\begin{itemize}
    \item \code{python -m pytest} for model and Streamlit state tests;
    \item \code{npm run test:browser} for mobile-touch and desktop-mouse
    Playwright workflows;
    \item \code{python scripts/benchmark\_fairness.py --budgets 1 2 3}
    for repeated fairness-phase stress trials at selected roster sizes; and
    \item \code{python scripts/benchmark\_hftb.py}\newline
    \code{--processes 4 --runs-per-process 5 --budget 5}\newline
    \code{--output benchmarks/hftb-5s-reference.json} for the opt-in 20-run HFTB latency,
    phase-timing, and lexicographic-quality contract across fresh processes.
\end{itemize}

The PDF was built with pdfTeX 3.141592653-2.6-1.40.29 (TeX Live 2026) by running
\code{pdflatex} twice with \code{-interaction=nonstopmode} and
\code{-halt-on-error}. Exact Python and OR-Tools versions are recorded for the
Team Red witness because the version ranges in \code{requirements.txt} are
deployment constraints rather than a fully pinned research environment.

\begingroup
\small
\setlength{\parskip}{2pt}
\begin{thebibliography}{9}

\bibitem{ortoolscpsat}
Google for Developers.
\href{https://developers.google.com/optimization/cp/cp_solver}
{\emph{CP-SAT Solver: OR-Tools}}. Accessed September 8, 2026.

\bibitem{ortoolssoftware}
Google OR-Tools.
\href{https://developers.google.com/optimization}
{\emph{OR-Tools Optimization Software}}. Accessed September 8, 2026.

\bibitem{streamlitdeploy}
Streamlit.
\href{https://docs.streamlit.io/deploy/streamlit-community-cloud/deploy-your-app/deploy}
{\emph{Deploy Your App on Community Cloud}}. Accessed September 8, 2026.

\bibitem{streamlitdeps}
Streamlit.
\href{https://docs.streamlit.io/deploy/streamlit-community-cloud/deploy-your-app/app-dependencies}
{\emph{App Dependencies for Your Community Cloud App}}. Accessed September 8,
2026.

\bibitem{bertsimas}
D. Bertsimas and J. N. Tsitsiklis.
\emph{Introduction to Linear Optimization}. Athena Scientific, 1997.

\bibitem{williams}
H. P. Williams.
\emph{Model Building in Mathematical Programming}, 5th ed. Wiley, 2013.

\end{thebibliography}
\endgroup

\end{document}
